\documentclass[12pt,a4paper]{article}
\usepackage[a4paper,margin=2.2cm]{geometry}
\usepackage[T2A]{fontenc}
\usepackage[utf8]{inputenc}
\usepackage[russian]{babel}
\usepackage{amsmath}
\usepackage{graphicx}
\usepackage{booktabs}
\usepackage{float}
\usepackage{hyperref}
\usepackage{array}
\usepackage{caption}
\usepackage{enumitem}
\captionsetup{font=small}
\setlist[itemize]{noitemsep, topsep=3pt}
\title{ТДЗ 18: VAR-модель региональной безработицы}
\author{Ковалёва Мария БЭК2310}
\date{}
\begin{document}
\maketitle
\thispagestyle{empty}
\section{Стационарность заданной VAR(1)}
Заданная система имеет вид:
\[
y_{1,t} = 1 + 0.6 y_{1,t-1} - 0.2 y_{2,t-1} + \varepsilon_{1,t},
\]
\[
y_{2,t} = 3 - 0.1 y_{1,t-1} + 0.4 y_{2,t-1} + \varepsilon_{2,t}.
\]
В матричной форме:
\[
Y_t = \alpha + B Y_{t-1} + \varepsilon_t,
\qquad
\alpha = \begin{pmatrix}1\\3\end{pmatrix},
\qquad
B = \begin{pmatrix}0.6 & -0.2\\-0.1 & 0.4\end{pmatrix}.
\]
Через лаговый оператор это записывается как
\[
(I - BL)Y_t = \alpha + \varepsilon_t.
\]
Стационарность проверяется по корням характеристического уравнения:
\[
\det(\lambda I - B) = \det\begin{pmatrix}\lambda - 0.6 & 0.2\\0.1 & \lambda - 0.4\end{pmatrix} = (\lambda - 0.6)(\lambda - 0.4) - 0.02 = \lambda^2 - \lambda + 0.22 = 0.
\]
Далее
\[
I-B = \begin{pmatrix}0.4 & 0.2\\0.1 & 0.6\end{pmatrix},
\qquad
(I-B)^{-1} = \begin{pmatrix}2.727273 & -0.909091\\-0.454545 & 1.818182\end{pmatrix}.
\]
И математическое ожидание:
\[
\mu = (I - B)^{-1}\alpha = \begin{pmatrix}2.727273 & -0.909091\\-0.454545 & 1.818182\end{pmatrix}\begin{pmatrix}1\\3\end{pmatrix}= \begin{pmatrix}0\\5\end{pmatrix}.
\]
\begin{table}[H]
\centering
\caption{Собственные значения матрицы $B$}
\label{tab:manual_stationarity}
\normalsize
\begin{tabular}{rrr}
\toprule
собственное значение & модуль & внутри единичной окружности \\
\midrule
0.6732 & 0.6732 & True \\
0.3268 & 0.3268 & True \\
\bottomrule
\end{tabular}
\end{table}
\begin{table}[H]
\centering
\caption{Безусловное математическое ожидание}
\label{tab:manual_mean}
\normalsize
\begin{tabular}{lr}
\toprule
компонента & математическое ожидание \\
\midrule
y1 & 0.0000 \\
y2 & 5.0000 \\
\bottomrule
\end{tabular}
\end{table}
Собственные значения равны 0.673205 и 0.326795; оба по модулю меньше единицы. Следовательно, VAR(1) стационарна и имеет конечное постоянное среднее $(0.000000, 5.000000)' = (0, 5)'$.

\section{Функция импульсного отклика}
Для VAR(1) ручные IRF определяются степенями матрицы перехода:
\[
IRF(1) = B = \begin{pmatrix}0.6 & -0.2\\-0.1 & 0.4\end{pmatrix},
\qquad
IRF(2) = B^2 = \begin{pmatrix}0.38 & -0.20\\-0.10 & 0.18\end{pmatrix}.
\]
\begin{table}[H]
\centering
\caption{Ручные значения IRF для заданной VAR(1)}
\label{tab:manual_irf}
\normalsize
\begin{tabular}{llrr}
\toprule
шок & отклик & t=1 & t=2 \\
\midrule
единичный шок в y1 & y1 & 0.6000 & 0.3800 \\
единичный шок в y1 & y2 & -0.1000 & -0.1000 \\
единичный шок в y2 & y1 & -0.2000 & -0.2000 \\
единичный шок в y2 & y2 & 0.4000 & 0.1800 \\
\bottomrule
\end{tabular}
\end{table}
\begin{figure}[H]
\centering
\includegraphics[width=0.78\textwidth]{figures/manual_var1_irf.png}
\caption{Схема ручных IRF для заданной VAR(1)}
\label{fig:manual_irf}
\end{figure}
Шок в $y_1$ дает положительный затухающий эффект на $y_1$ и слабый отрицательный эффект на $y_2$. Шок в $y_2$ дает отрицательный эффект на $y_1$ и положительный затухающий эффект на $y_2$. Поскольку система стационарна, отклики по мере роста горизонта уменьшаются по модулю.

\section{Данные и выбор регионов}
\begin{itemize}
\item переменная: уровень безработицы, в процентах;
\item частота: годовая;
\item число наблюдений в уровнях: 23;
\item модель далее строится не по уровням, а по первым разностям.
\end{itemize}
\begin{figure}[H]
\centering
\includegraphics[width=0.92\textwidth]{figures/unemployment_selected_regions_combined.png}
\caption{Исходные ряды уровней безработицы}
\label{fig:data}
\end{figure}
Используются годовые данные по уровню безработицы за 2000--2022 гг. Выбраны три соседних региона: Белгородская, Курская и Воронежская области. Такой выбор позволяет сохранить географическую и экономическую сопоставимость рядов.

\section{Предварительный анализ}
На графиках видно общее снижение безработицы с отдельными всплесками. Следовательно, в уровнях ряды выглядят не вполне стационарными, и это нужно проверить формально. Дополнительная оговорка: после перехода к первым разностям остается около 22 наблюдений, то есть выборка мала.

\section{Стационарность рядов}
\begin{table}[H]
\centering
\caption{Тесты стационарности}
\label{tab:stationarity}
\resizebox{\textwidth}{!}{%
\begin{tabular}{llrrl}
\toprule
регион & форма\_ряда & ADF p-value & KPSS p-value & итог \\
\midrule
Белгородская область & Уровни & 0.1515 & 0.0433 & Ряд нестационарен \\
Курская область & Уровни & 0.3388 & 0.0184 & Ряд нестационарен \\
Воронежская область & Уровни & 0.4156 & 0.0197 & Ряд нестационарен \\
Белгородская область & Первые разности & 0.0018 & 0.1000 & Ряд стационарен \\
Курская область & Первые разности & 0.0001 & 0.1000 & Ряд стационарен \\
Воронежская область & Первые разности & 0.0013 & 0.1000 & Ряд стационарен \\
\bottomrule
\end{tabular}
}
\end{table}
По уровням ADF не отвергает единичный корень, а KPSS отвергает стационарность. По первым разностям картина становится обратной: ряды можно считать стационарными. Поэтому итоговая VAR оценивается на первых разностях.

\section{Выбор лага и спецификация VAR}
Эмпирическая модель записывается как
\[
\Delta Y_t = c + A_1 \Delta Y_{t-1} + A_2 \Delta Y_{t-2} + A_3 \Delta Y_{t-3} + \varepsilon_t.
\]
\begin{table}[H]
\centering
\caption{Выбор длины лага}
\label{tab:lags}
\resizebox{\textwidth}{!}{%
\begin{tabular}{rrrrrrr}
\toprule
лаг & эффективные наблюдения & параметров на уравнение & AIC & BIC & HQIC & FPE \\
\midrule
1 & 21 & 4 & 0.0695 & 0.6664 & 0.1990 & 1.0872 \\
2 & 20 & 7 & -1.6353 & -0.5898 & -1.4312 & 0.2138 \\
3 & 19 & 10 & -2.2163 & -0.7251 & -1.9640 & 0.1550 \\
\bottomrule
\end{tabular}
}
\end{table}
\begin{table}[H]
\centering
\caption{Небольшая проверка устойчивости вывода по лагам}
\label{tab:robustness}
\resizebox{\textwidth}{!}{%
\begin{tabular}{rrrrrl}
\toprule
лаг & VAR устойчива & p-value Portmanteau & p-value Portmanteau adjusted & p-value нормальности & значимые Granger-связи для изменений безработицы \\
\midrule
1 & True & 0.0944 & 0.0164 & 0.0000 & нет \\
2 & True & 0.0209 & 0.0023 & 0.8150 & Воронежская область -> Курская область \\
3 & True & 0.1656 & 0.0670 & 0.1890 & Воронежская область -> Курская область \\
\bottomrule
\end{tabular}
}
\end{table}
Информационные критерии AIC, BIC и HQIC выбирают лаг $p = 3$. Итоговая спецификация --- \textbf{VAR(3) на первых разностях уровней безработицы}. Это решение поддерживается диагностикой и сравнением с VAR(1) и VAR(2), но из-за малой выборки VAR(3) нужно трактовать осторожно.

\section{Диагностика модели}
\begin{table}[H]
\centering
\caption{Диагностика остатков VAR(3)}
\label{tab:diag}
\resizebox{\textwidth}{!}{%
\begin{tabular}{lrl}
\toprule
тест & p-value & вывод \\
\midrule
Portmanteau на автокорреляцию остатков & 0.1656 & H0 не отвергается \\
Portmanteau на автокорреляцию остатков (small-sample adjusted) & 0.0670 & H0 не отвергается \\
Тест нормальности остатков & 0.1890 & H0 не отвергается \\
\bottomrule
\end{tabular}
}
\end{table}
\begin{figure}[H]
\centering
\includegraphics[width=0.60\textwidth]{figures/var_inverse_roots.png}
\caption{Собственные значения companion-матрицы}
\label{fig:roots}
\end{figure}
VAR(3) устойчива: собственные значения лежат внутри единичной окружности. Тесты на автокорреляцию и нормальность остатков на уровне 5\% не дают критических возражений. Для учебной задачи модель можно считать приемлемой, хотя малая выборка остается ограничением.

\section{Причинность по Грейнджеру}
\begin{table}[H]
\centering
\caption{Тесты Грейнджера}
\label{tab:granger}
\resizebox{\textwidth}{!}{%
\begin{tabular}{lrl}
\toprule
направление & p-value & вывод \\
\midrule
Воронежская область -> Курская область & 0.0068 & Отвергаем H0: изменения безработицы в X помогают прогнозировать изменения безработицы в Y \\
Курская область -> Воронежская область & 0.1696 & Не отвергаем H0: значимой Granger-связи для изменений безработицы не найдено \\
Белгородская область -> Воронежская область & 0.3372 & Не отвергаем H0: значимой Granger-связи для изменений безработицы не найдено \\
Белгородская область -> Курская область & 0.3890 & Не отвергаем H0: значимой Granger-связи для изменений безработицы не найдено \\
Воронежская область -> Белгородская область & 0.5661 & Не отвергаем H0: значимой Granger-связи для изменений безработицы не найдено \\
Курская область -> Белгородская область & 0.8380 & Не отвергаем H0: значимой Granger-связи для изменений безработицы не найдено \\
\bottomrule
\end{tabular}
}
\end{table}
Так как модель оценена на первых разностях, выводы относятся к \textbf{изменениям безработицы}. Единственная статистически значимая связь: \textbf{Воронежская область -> Курская область}. Это предиктивная, а не структурная причинность.

\section{Функции импульсного отклика}
\begin{figure}[H]
\centering
\includegraphics[width=0.92\textwidth]{figures/var_irf_orth_5y.png}
\caption{IRF для VAR(3) на первых разностях}
\label{fig:irf}
\end{figure}
IRF интерпретируются для \textbf{изменений безработицы}. Наиболее заметная межрегиональная реакция наблюдается для отклика Курской области на шок в Воронежской области. Знаки откликов могут меняться, но на дальних горизонтах реакции затухают, что согласуется с устойчивостью VAR.

\section{Разложение дисперсии ошибки прогноза}
\begin{table}[H]
\centering
\caption{FEVD на горизонте 5 лет}
\label{tab:fevd5}
\resizebox{\textwidth}{!}{%
\begin{tabular}{lrrrr}
\toprule
отклик & горизонт & Белгородская область & Курская область & Воронежская область \\
\midrule
Белгородская область & 5 & 0.8674 & 0.0070 & 0.1255 \\
Курская область & 5 & 0.4095 & 0.1430 & 0.4475 \\
Воронежская область & 5 & 0.4623 & 0.0729 & 0.4648 \\
\bottomrule
\end{tabular}
}
\end{table}
\begin{figure}[H]
\centering
\includegraphics[width=0.92\textwidth]{figures/var_fevd_5y.png}
\caption{Разложение дисперсии ошибки прогноза}
\label{fig:fevd}
\end{figure}
FEVD относится к \textbf{изменениям безработицы}. Для Белгородской области доминируют собственные шоки (0.867). Для Курской области вклад внешних шоков велик: Воронежская область дает 0.448, Белгородская --- 0.409, что вместе больше собственной доли 0.143. Для Воронежской области собственная доля 0.465, но влияние Белгородской области тоже заметно (0.462). Значит, межрегиональная динамика особенно важна для Курской области.

\section{Прогноз}
\begin{table}[H]
\centering
\caption{Прогноз изменений безработицы}
\label{tab:forecast_changes}
\resizebox{\textwidth}{!}{%
\begin{tabular}{rrrr}
\toprule
год & изменение: Белгородская область & изменение: Курская область & изменение: Воронежская область \\
\midrule
2023 & 0.1459 & -0.2861 & -1.3666 \\
2024 & -0.3644 & -1.3242 & -0.7130 \\
2025 & -0.3031 & -0.5164 & -0.4661 \\
\bottomrule
\end{tabular}
}
\end{table}
\begin{table}[H]
\centering
\caption{Приближенный прогноз уровней безработицы}
\label{tab:forecast_levels}
\resizebox{\textwidth}{!}{%
\begin{tabular}{rrrr}
\toprule
год & уровень: Белгородская область & уровень: Курская область & уровень: Воронежская область \\
\midrule
2023 & 3.8459 & 2.6139 & 2.1334 \\
2024 & 3.4815 & 1.2897 & 1.4204 \\
2025 & 3.1784 & 0.7734 & 0.9543 \\
\bottomrule
\end{tabular}
}
\end{table}
\begin{figure}[H]
\centering
\includegraphics[width=0.92\textwidth]{figures/var_forecast_changes_levels.png}
\caption{Краткосрочный прогноз по VAR(3)}
\label{fig:forecast}
\end{figure}
Сначала прогнозируются изменения безработицы, а затем уровни восстанавливаются механически от значения 2022 года. Из 9 прогнозных значений изменений безработицы 8 отрицательны, то есть модель чаще ожидает дальнейшее снижение безработицы. Однако прогноз следует считать только краткосрочным и исследовательским из-за малого объема выборки.

\section{Итоговый вывод}
Заданная VAR(1) является стационарной, ее импульсные отклики затухают, а безусловное среднее конечно. В эмпирической части модель \textbf{VAR(3) на первых разностях} для Белгородской, Курской и Воронежской областей показывает устойчивую, но достаточно чувствительную к малой выборке динамику. Наиболее содержательный результат --- связь от изменений безработицы в Воронежской области к изменениям безработицы в Курской области; это подтверждается и тестами Грейнджера, и IRF, и FEVD. Следовательно, динамика безработицы в выбранных регионах частично взаимосвязана, но количественные выводы и прогноз следует трактовать осторожно из-за малого числа наблюдений.

\end{document}
